%%%%%%%%%%%%%%%%%%%%%%%%%%%%% Define Article %%%%%%%%%%%%%%%%%%%%%%%%%%%%%%%%%%
\documentclass[11pt, a4paper]{article}
%%%%%%%%%%%%%%%%%%%%%%%%%%%%%%%%%%%%%%%%%%%%%%%%%%%%%%%%%%%%%%%%%%%%%%%%%%%%%%%

%%%%%%%%%%%%%%%%%%%%%%%%%%%%% Using Packages %%%%%%%%%%%%%%%%%%%%%%%%%%%%%%%%%%
\usepackage{geometry}
\usepackage{graphicx}
\usepackage{amssymb}
\usepackage{amsmath}
\usepackage{amsthm}
\usepackage{empheq}
\usepackage{mdframed}
\usepackage{booktabs}
\usepackage{lipsum}
\usepackage{color}
\usepackage{psfrag}
\usepackage{pgfplots}
\usepackage{bm}
\usepackage{float}
\usepackage[english]{babel}
\usepackage{csquotes}
\usepackage{biblatex}
\addbibresource{refs.bib}
\usepackage{authblk}
%%%%%%%%%%%%%%%%%%%%%%%%%%%%%%%%%%%%%%%%%%%%%%%%%%%%%%%%%%%%%%%%%%%%%%%%%%%%%%%

% Other Settings

%%%%%%%%%%%%%%%%%%%%%%%%%% Page Setting %%%%%%%%%%%%%%%%%%%%%%%%%%%%%%%%%%%%%%%
\geometry{a4paper, margin=1in}
\graphicspath{{img/}}
%%%%%%%%%%%%%%%%%%%%%%%%%% Define some useful colors %%%%%%%%%%%%%%%%%%%%%%%%%%
\definecolor{ocre}{RGB}{243,102,25}
\definecolor{mygray}{RGB}{243,243,244}
\definecolor{deepGreen}{RGB}{26,111,0}
\definecolor{shallowGreen}{RGB}{235,255,255}
\definecolor{deepBlue}{RGB}{61,124,222}
\definecolor{shallowBlue}{RGB}{235,249,255}
%%%%%%%%%%%%%%%%%%%%%%%%%%%%%%%%%%%%%%%%%%%%%%%%%%%%%%%%%%%%%%%%%%%%%%%%%%%%%%%

%%%%%%%%%%%%%%%%%%%%%%%%%% Define an orangebox command %%%%%%%%%%%%%%%%%%%%%%%%
\newcommand\orangebox[1]{\fcolorbox{ocre}{mygray}{\hspace{1em}#1\hspace{1em}}}
%%%%%%%%%%%%%%%%%%%%%%%%%%%%%%%%%%%%%%%%%%%%%%%%%%%%%%%%%%%%%%%%%%%%%%%%%%%%%%%

%%%%%%%%%%%%%%%%%%%%%%%%%%%% English Environments %%%%%%%%%%%%%%%%%%%%%%%%%%%%%
\newtheoremstyle{mytheoremstyle}{3pt}{3pt}{\normalfont}{0cm}{\rmfamily\bfseries}{}{1em}{{\color{black}\thmname{#1}~\thmnumber{#2}}\thmnote{\,--\,#3}}
\newtheoremstyle{myproblemstyle}{3pt}{3pt}{\normalfont}{0cm}{\rmfamily\bfseries}{}{1em}{{\color{black}\thmname{#1}~\thmnumber{#2}}\thmnote{\,--\,#3}}
\theoremstyle{mytheoremstyle}
\newmdtheoremenv[linewidth=1pt,backgroundcolor=shallowGreen,linecolor=deepGreen,leftmargin=0pt,innerleftmargin=20pt,innerrightmargin=20pt,]{theorem}{Theorem}[section]
\theoremstyle{mytheoremstyle}
\newmdtheoremenv[linewidth=1pt,backgroundcolor=shallowBlue,linecolor=deepBlue,leftmargin=0pt,innerleftmargin=20pt,innerrightmargin=20pt,]{definition}{Definition}[section]
\theoremstyle{myproblemstyle}
\newmdtheoremenv[linecolor=black,leftmargin=0pt,innerleftmargin=10pt,innerrightmargin=10pt,]{problem}{Problem}[section]
%%%%%%%%%%%%%%%%%%%%%%%%%%%%%%%%%%%%%%%%%%%%%%%%%%%%%%%%%%%%%%%%%%%%%%%%%%%%%%%

%%%%%%%%%%%%%%%%%%%%%%%%%%%%%%% Plotting Settings %%%%%%%%%%%%%%%%%%%%%%%%%%%%%
\usepgfplotslibrary{colorbrewer}
\pgfplotsset{width=8cm,compat=1.9}
%%%%%%%%%%%%%%%%%%%%%%%%%%%%%%%%%%%%%%%%%%%%%%%%%%%%%%%%%%%%%%%%%%%%%%%%%%%%%%%

%%%%%%%%%%%%%%%%%%%%%%%%%%%%%%% Title & Author %%%%%%%%%%%%%%%%%%%%%%%%%%%%%%%%
\title{Team ESCOM-IPN at CLEF 2026: Multimodal Financial Decision Making using LLMs, Chronos, and Auto-Reflective Prompt Routing}
\author[1]{Rodrigo}
\author[1]{Jona}
\author[1]{Baru}
\author[1]{David}
\affil[1]{Escuela Superior de Cómputo (ESCOM) - Instituto Politécnico Nacional (IPN), Mexico City, Mexico}
\date{}
%%%%%%%%%%%%%%%%%%%%%%%%%%%%%%%%%%%%%%%%%%%%%%%%%%%%%%%%%%%%%%%%%%%%%%%%%%%%%%%

\begin{document}
    \maketitle
    
    \begin{abstract}
    This Working Note describes the participation of the ESCOM-IPN team in Task 3 of the CLEF 2026 FinMMEval lab. We propose a highly structured, multimodal retrieval-augmented generation (RAG) pipeline to predict financial trading actions (BUY, HOLD, SELL) across diverse assets. Our system integrates unstructured news and SEC filings with numerical scalars and temporal price sequences. We introduce a novel dual-strategy \textit{PromptBuilder} featuring a tripartite temporal memory (rolling buffer, self-reflection, and top-K retrieval), allowing Large Language Models (LLMs) to reason over current market contexts while continuously learning from empirical trading results. Furthermore, we implement \textit{finvert}, a highly optimized FinBERT-based module for robust market sentiment extraction. Our methodology demonstrates how programmatic prompt engineering, heuristic risk management, and chronological FIFO queuing of semantic summaries can significantly enhance LLM reasoning in simulated financial environments.
    \end{abstract}

    \textbf{Keywords:} Financial Forecasting, Large Language Models, Sentiment Analysis, Prompt Engineering, Temporal Memory, FinMMEval, CLEF 2026.

    \section{Introduction}
    Financial decision-making requires the rapid synthesis of heterogeneous data sources. In the CLEF 2026 FinMMEval Lab (Task 3), participants are challenged to develop agents capable of processing numerical scalars, temporal price histories, and textual news to recommend trading actions (BUY, HOLD, SELL) for assets such as Bitcoin, Tesla, and Apple. 

    While Large Language Models (LLMs) excel at natural language understanding, they traditionally struggle with raw numerical sequences and lack inherent contextual memory of past actions in simulated trading environments. To bridge this gap, our team developed a comprehensive multimodal pipeline. 

    The core contribution of our work lies in the intelligent orchestration of data prior to LLM inference. Specifically, we detail two critical components designed by our team: (1) \textbf{finvert}, an automated and resource-optimized sentiment extraction module leveraging FinBERT, and (2) the \textbf{PromptBuilder}, a sophisticated prompt engineering engine equipped with a tripartite temporal memory. By providing the LLM not just with current data, but with a rolling history of its past decisions and their resulting financial rewards, we simulate a realistic trading agent that adapts to market dynamics and mitigates risk through self-reflection.

    \section{System Architecture and Dataset Formulation}
    \textit{[Notau: Expandir esta sección. Describir script'generate\_all.py', cómo se construyó 'finmmeval\_raw.csv' mezclando HuggingFace y Yahoo Finance, y los detalles de Chronos Encoder.]}

    The overall architecture transforms raw multimodal data into structured prompts for LLM inference. The pipeline operates as follows:
    \begin{enumerate}
        \item \textbf{Data Ingestion:} Raw data from CLEF and Yahoo Finance is collected and unified, encompassing textual news, 10-K/10-Q filings, and technical scalars.
        \item \textbf{Temporal Encoding (Chronos):} Time-series data is processed to extract temporal embeddings. \textit{[David: Add Chronos details here]}.
        \item \textbf{Sentiment Extraction (finvert):} Textual data is classified to extract sentiment probabilities.
        \item \textbf{Prompt Orchestration:} The \textit{PromptBuilder} fuses all signals alongside the agent's trading memory.
        \item \textbf{LLM Decision Making:} The final prompt is submitted to the decision LLM to output a structured JSON containing the action and rationale.
    \end{enumerate}

    \section{Dataset and Feature Engineering}
    \label{sec:dataset}

    \subsection{Data sources}
    \label{sec:datasources}

    The primary source is the \textit{TheFinAI/daily\_news} dataset~\cite{theFinAI2025}, publicly available through HuggingFace, comprising daily financial records for twelve assets over approximately 268 trading days, from August 1, 2025 to April 25, 2026. The covered assets are ten listed equities (TSLA, MSFT, AAPL, AMZN, GOOGL, META, NVDA, ADBE, BMRN, and MRNA) and two cryptocurrencies (BTC and ETH).

    Each daily record contains the closing price, a news summary of between 800 and 1\,500 words, SEC regulatory filings (forms 10-K and 10-Q) when available, a momentum label (bullish, neutral, or bearish), and the next-day price difference --- used exclusively as a reward signal during training and never exposed to the model as an input.

    As a complementary source, historical prices from Yahoo Finance were used for a four-month warm-up period prior to the start of the dataset (April--July 2025). Their sole purpose was to ensure that indicators with long lookback windows --- the 60-day moving average regime signal and the 30-day rolling Sharpe ratio --- could be computed without producing missing values at the start of the main period. These data were not incorporated as training rows, as doing so would have created an imbalance across assets since only TSLA and BTC had extended historical data of this type.

    \subsection{Feature engineering}
    \label{sec:features}

    A total of 43 candidate features were constructed across six categories. The common design criterion is the absence of temporal information leakage: every feature at time~$t$ uses only information available up to that point.

    \paragraph{Returns and trend.}
    Log returns are computed at one, three, and five days. The price position relative to its 10-, 20-, and 50-day moving averages is expressed as a percentage deviation. A binary market regime signal takes value $+1$ when the price exceeds the 60-day moving average (bullish regime) and $-1$ otherwise.

    \paragraph{Volatility.}
    Annualized realized volatility is computed over windows of 5, 10, 20, and 30 days. Two volatility ratios --- $\text{vol}_{5d}/\text{vol}_{20d}$ and $\text{vol}_{10d}/\text{vol}_{30d}$ --- detect the expansion or contraction of short-term volatility relative to longer baseline levels, a signal particularly relevant for cryptocurrencies. A binary volatility regime flag activates when the 10-day volatility exceeds the historical 70th percentile.

    \paragraph{Technical oscillators and MACD.}
    The RSI is computed at periods of 9 and 14 days and normalized to $[0, 1]$. Bollinger Bands contribute the percentage position of the price within the bands and the relative band width. MACD is decomposed into line, signal, and histogram components; the histogram is particularly useful for detecting imminent trend crossovers.

    \paragraph{Agent internal state.}
    The rolling Sharpe ratio at 30 and 10 days reflects the risk-adjusted performance of recent decisions. State confidence accumulates consecutive trading days without a significant price shock (return exceeding $3\sigma$), normalized to $[0, 1]$, following the FINMEM framework~\cite{finmem2024}.

    \paragraph{Text-derived features.}
    Sentiment analysis is performed with ProsusAI/FinBERT~\cite{araci2019finbert}, a model specialized in financial texts. Since news summaries frequently exceed the 512-token limit, each document is split into 450-token chunks. Scores are averaged across chunks with a weight of $1.5\times$ applied to the first, as the introductory paragraph tends to concentrate the dominant sentiment signal. For regulatory filings (10-K and 10-Q), the same chunked procedure is applied; since filings appear at most once per quarter, the computed score is propagated forward using last observation carried forward (LOCF), and the number of days since the last available filing --- normalized by 90 --- is included to model the temporal decay of filing relevance.

    \subsection{Feature selection}
    \label{sec:selection}

    Selection was performed jointly across all twelve assets (approximately 3\,200 observations), providing more reliable mutual information estimates than per-asset analysis. The process consists of three sequential steps.

    \textbf{(1) Spearman correlation filter.} Feature pairs with an absolute Spearman rank correlation coefficient exceeding 0.85 were identified. Spearman was preferred over Pearson because financial features exhibit asymmetric distributions and non-linear relationships. From each highly correlated pair, the feature with lower mutual information with respect to the target was discarded.

    \textbf{(2) Mutual information filter.} Mutual information between each feature and the binary target --- next-day price direction --- was estimated using the $k$-nearest neighbors method with $k=5$~\cite{kraskov2004}. Features scoring below 0.005 were discarded.

    \textbf{(3) Temporal stability.} A two-sample Kolmogorov-Smirnov test was applied between the feature distributions of the first and last temporal folds. Features with KS statistic exceeding 0.30 were flagged for review, as distributional drift in financial series may reflect genuine structural regime changes rather than data quality issues. The analysis was run independently for equities and cryptocurrencies.

    \subsection{Final feature sets}
    \label{sec:finalfeatures}

    After selection, 16 features were retained for equities and 16 for cryptocurrencies. Tables~\ref{tab:features_equity} and~\ref{tab:features_crypto} detail the final sets. Both share a core of ten common features --- returns, news sentiment, encoded momentum, and basic metadata --- and differ in the technical signals with the highest predictive power for each asset class.

    \begin{table}[H]
    \centering
    \begin{tabular}{llll}
    \toprule
    \textbf{Feature} & \textbf{Range} & \textbf{MI} & \textbf{Rationale} \\
    \midrule
    \texttt{asset\_type}           & $\{0, 1\}$       & ---    & Asset class identifier \\
    \texttt{day\_of\_week}         & $[0, 6]$         & ---    & Calendar effects \\
    \texttt{month}                 & $[1, 12]$        & ---    & Monthly seasonality \\
    \texttt{has\_news}             & $\{0, 1\}$       & ---    & News availability \\
    \texttt{has\_momentum}         & $\{0, 0.5, 1\}$  & ---    & Momentum data quality \\
    \texttt{finbert\_news\_full}   & $[-1, 1]$        & ---    & Sentiment (full text) \\
    \texttt{finbert\_news\_intro}  & $[-1, 1]$        & ---    & Sentiment (intro paragraph) \\
    \texttt{momentum\_encoded}     & $\{-1, 0, 1\}$   & ---    & Momentum direction \\
    \texttt{log\_return\_1d}       & $[-0.5, 0.5]$    & ---    & Base directional signal \\
    \texttt{log\_return\_3d}       & $[-0.8, 0.8]$    & ---    & Short-term trend \\
    \texttt{state\_confidence}     & $[0, 1]$         & Med.   & Context reliability \\
    \texttt{volatility\_20d}       & $[0, 3]$         & Med.   & Asset risk level \\
    \texttt{rsi\_9}                & $[0, 1]$         & High   & Overbought/oversold (fast) \\
    \texttt{rsi\_14}               & $[0, 1]$         & High   & Overbought/oversold (standard) \\
    \texttt{news\_volume\_proxy}   & $[0, 1]$         & Med.   & Market attention \\
    \texttt{momentum\_strength}    & $[0, 5]$         & Med.   & Vol.-adjusted movement strength \\
    \bottomrule
    \end{tabular}
    \caption{Final selected features --- Equities.}
    \label{tab:features_equity}
    \end{table}

    \begin{table}[H]
    \centering
    \begin{tabular}{llll}
    \toprule
    \textbf{Feature} & \textbf{Range} & \textbf{MI} & \textbf{Rationale} \\
    \midrule
    \texttt{asset\_type}           & $\{0, 1\}$       & ---    & Asset class identifier \\
    \texttt{day\_of\_week}         & $[0, 6]$         & ---    & Calendar effects (incl. weekends) \\
    \texttt{month}                 & $[1, 12]$        & ---    & Monthly seasonality \\
    \texttt{has\_news}             & $\{0, 1\}$       & ---    & News availability \\
    \texttt{has\_momentum}         & $\{0, 0.5, 1\}$  & ---    & Momentum data quality \\
    \texttt{finbert\_news\_full}   & $[-1, 1]$        & ---    & Sentiment (full text) \\
    \texttt{finbert\_news\_intro}  & $[-1, 1]$        & ---    & Sentiment (intro paragraph) \\
    \texttt{momentum\_encoded}     & $\{-1, 0, 1\}$   & ---    & Momentum direction \\
    \texttt{log\_return\_1d}       & $[-0.5, 0.5]$    & ---    & Base directional signal \\
    \texttt{log\_return\_3d}       & $[-0.8, 0.8]$    & ---    & Short-term trend \\
    \texttt{sr\_30d}               & $[-3, 3]$        & High   & Risk-adjusted recent performance \\
    \texttt{price\_vs\_ma\_10}     & $[-30, 30]\%$    & High   & Short-term relative position \\
    \texttt{rsi\_14}               & $[0, 1]$         & High   & Overbought/oversold \\
    \texttt{vol\_ratio\_5\_20}     & $[0.3, 4.0]$     & Med.   & Volatility expansion \\
    \texttt{state\_confidence}     & $[0, 1]$         & Med.   & Context reliability \\
    \texttt{vol\_ratio\_10\_30}    & $[0.3, 4.0]$     & Med.   & Volatility regime \\
    \bottomrule
    \end{tabular}
    \caption{Final selected features --- Cryptocurrencies.}
    \label{tab:features_crypto}
    \end{table}

    \subsection{Text treatment}
    \label{sec:text}

    News summaries and regulatory filings are stored separately from the numerical features and passed directly to the LLM as natural language. This decision rests on three observations.

    First, FinBERT sentiment scores obtained near-zero mutual information values in both asset classes. The dataset's texts are pre-processed summaries that frequently contain sentiment-neutralizing language --- phrases such as \textit{``overall market sentiment remains neutral''} cause FinBERT to underestimate the actual directional signal. Compressing 800--1\,500 words into a single scalar discards the causal and entity-specific information that gives the text its predictive value. Second, large language models are significantly more capable of extracting actionable signals from structured financial prose than scalar sentiment classifiers~\cite{zhang2023fingpt}. Third, keeping the text separate preserves flexibility with respect to the target model's context window.

    FinBERT scores are nonetheless retained in the numerical feature set as supplementary signals, not as the primary sentiment input.

    % References are in refs.bib — loaded via \addbibresource in preamble

    \section{Temporal Encoding: The Chronos-Bolt Module}
    To capture the predictive signal embedded in historical price trajectories, our pipeline employs \textbf{Chronos-Bolt} (\texttt{amazon/chronos-bolt-small}), a transformer-based probabilistic forecasting model pre-trained on large collections of real-world time series. Unlike traditional statistical approaches, Chronos-Bolt learns a distribution over future values, enabling the extraction of uncertainty-aware forecasts directly from raw numerical sequences.

    The module is loaded via the \texttt{BaseChronosPipeline} abstraction, which handles device placement automatically. When a CUDA-capable GPU is available, the model is instantiated in \texttt{bfloat16} precision to reduce VRAM consumption while preserving numerical stability. In CPU-only environments, standard floating-point inference is applied as a fallback.

    \subsection{Input Formulation and Batch Processing}
    Each asset's historical prices are stored per-day as a serialized list in the \texttt{prices\_sequence} column of the dataset. Prior to inference, each entry is deserialized into a typed list of floats, forming the context window fed to the model.

    To maximize GPU throughput, sequences are grouped into batches of 16 and converted to PyTorch tensors. For each batch, the model is queried with \texttt{prediction\_length=1}, restricting the forecast horizon to a single future time step — the next trading day. This deliberate constraint enforces that the model can only reason over observed history, preserving temporal isolation and avoiding any form of look-ahead bias.

    \subsection{Forecast Extraction and Direction Classification}
    Chronos-Bolt outputs a probabilistic forecast tensor of shape $(\text{batch} \times \text{num\_samples} \times \text{prediction\_length})$, where the \texttt{num\_samples} dimension represents draws from the learned predictive distribution. To obtain a single point estimate per sequence, we compute the \textbf{median} across the sample dimension:

    \begin{equation}
        \hat{p}_{t+1} = \mathrm{median}_{s} \; F_s(x_{1:t})
    \end{equation}

    where $F_s$ denotes the $s$-th sample of the forecast distribution conditioned on the observed price history $x_{1:t}$. The resulting scalar $\hat{p}_{t+1}$ is then converted into a discrete directional class by computing the relative price change with respect to the last observed price $p_t$:

    \begin{equation}
        \Delta = \frac{\hat{p}_{t+1} - p_t}{p_t}
    \end{equation}

    A symmetric threshold $\tau = 0.001$ ($\pm$0.1\%) governs the classification into three mutually exclusive categories:

    \begin{equation}
        \text{class} =
        \begin{cases}
            1 & \text{if } \Delta > \tau \quad (\textsc{up}) \\
            2 & \text{if } \Delta < -\tau \quad (\textsc{down}) \\
            0 & \text{otherwise} \quad (\textsc{hold})
        \end{cases}
    \end{equation}

    This threshold was selected to filter out negligible price fluctuations below the level of market noise, ensuring that the \textsc{hold} class captures genuine price stability rather than prediction imprecision.

    \subsection{Evaluation and Output}
    Upon completing inference over the full dataset, the module computes both global and per-asset evaluation metrics. Global performance is assessed via accuracy, macro-averaged F1, and weighted F1 scores, while per-asset metrics allow identification of assets whose temporal price patterns are more or less amenable to the model's forecasting strategy. The complete predictions and metrics are persisted to \texttt{predictions.csv} and \texttt{metrics.csv} respectively, serving as structured inputs to subsequent stages of the pipeline.

    \section{Multimodal Sentiment Extraction: The \textit{finvert} Module}
    The \textit{finvert} module acts as the market sentiment sensor of our architecture. Its primary objective is to process unstructured financial texts and quantify them into actionable numerical probabilities.

    \subsection{Auto-Discovery and Text Combination}
    Financial sentiment is rarely derived from a single source. Our system aggregates three primary textual streams per day in a strict priority hierarchy: daily news, annual reports (10-K), and quarterly reports (10-Q). 
    
    To optimize computational resources, the module implements an auto-discovery and filtering mechanism (\texttt{has\_real\_text}). Datasets lacking actual textual data (e.g., pure price-action assets from Yahoo Finance flagged as \texttt{[NO\_TEXT]}) are automatically bypassed. These rows are deterministically assigned a \textit{neutral} label with a 1.0 confidence score without invoking the GPU, vastly improving pipeline efficiency and preventing unnecessary VRAM allocation.

    \subsection{FinBERT Inference and Batch Processing}
    For valid textual entries, we utilize \textit{FinBERT} (ProsusAI/finbert), a BERT model pre-trained on large-scale financial corpora. Texts are processed in batches of 32 to maximize GPU throughput. 

    The module tokenizes the combined daily text (truncated to a maximum length of 512 tokens) and outputs logits through a 3-class classification head. Applying a softmax function, we extract the probabilities for each class:
    
    \begin{equation}
        P(C_i | x) = \frac{e^{z_i}}{\sum_{j=1}^{3} e^{z_j}}
    \end{equation}
    
    where $C \in \{\text{positive}, \text{negative}, \text{neutral}\}$. The output is stored in a structured format containing the \texttt{sentiment\_label}, the absolute \texttt{sentiment\_score}, and the exact probability distribution. This granular output prevents the downstream LLM from interpreting a weak positive signal with the same weight as a strong one.

    \section{Dynamic Prompt Routing and Temporal Memory}
    The \textit{PromptBuilder} is the central nervous system of our pipeline, encompassing over 600 lines of orchestration logic. Its responsibility is to synthesize sentiment analysis, technical scalars, portfolio status, and temporal sequences into a highly structured prompt. 
    
    Crucially, to simulate a realistic trading environment and prevent look-ahead bias, we enforce \textbf{strict temporal isolation}: the prompt for day $T$ is generated entirely \textit{before} the portfolio memory is updated with the market outcomes of day $T$.

    \subsection{Tripartite Memory Architecture}
    At any time step $t$, the state of the market is captured by a \texttt{SensorOutputs} object. Meanwhile, the system maintains a \texttt{MemoryContext} of \texttt{DayRecord} objects. Inspired by recent advancements in LLM trading agents (such as FLAG-TRADER and FINMEM), we implemented three types of memory:
    \begin{itemize}
        \item \textbf{Rolling Buffer:} A sequential tracker of recent actions and rewards.
        \item \textbf{Self-Reflection:} Analysis of past errors and successes.
        \item \textbf{Top-K Retrieval:} Extraction of the most relevant historical trades.
    \end{itemize}

    \subsection{Strategy V1: Programmatic Auto-Reflection}
    Strategy V1 is a deterministic, highly efficient prompt routing mechanism requiring no auxiliary API calls. It programmatically translates numerical scalars (e.g., RSI, portfolio exposure) and categorical data into natural language. 

    The primary innovation in V1 is the \textbf{Auto-Reflection Summary}. The module identifies the most successful and most detrimental trades within a historical window. The generated prompt presents the LLM with an ASCII-formatted rolling history table and a specific reflection block:
    
    \begin{mdframed}[linewidth=1pt,backgroundcolor=mygray,linecolor=deepBlue]
    \textit{"CASOS EXITOSOS RECIENTES: 2024-01-15: Accion=BUY, Reward=+0.0234... \\
    CASOS NEGATIVOS RECIENTES: 2024-01-17: Accion=SELL, Reward=-0.0120..."}
    \end{mdframed}
    
    This structural constraint forces the LLM to ground its current decision by learning from empirical past mistakes, effectively mimicking human reinforcement learning.

    \subsection{Strategy V2: Dual-LLM Pipeline and Semantic FIFO Queuing}
    To further enhance reasoning, Strategy V2 employs a chained dual-LLM architecture, operating as follows:
    \begin{enumerate}
        \item \textbf{LLM-1 (Summarizer):} A lightweight model (e.g., GPT-4o-mini or Llama-3.2-3B) reads the raw data of past days and translates the events into a rich, narrative summary. To ensure deterministic testing during development, we also implemented a \texttt{StubClient} factory pattern.
        \item \textbf{FIFO Queuing \& Caching:} These narrative summaries are pushed into a chronological First-In-First-Out (FIFO) queue. To heavily optimize API costs and reduce latency, outputs from LLM-1 are cached directly in the \texttt{DayRecord.llm\_summary} attribute.
        \item \textbf{LLM-2 (Decision):} A heavier reasoning model (e.g., GPT-4o) receives the entire narrative FIFO queue alongside the current market context to make the final BUY/HOLD/SELL decision.
    \end{enumerate}

    \subsection{Heuristic Risk Management via System Prompts}
    Both strategies are governed by a strict \textit{System Prompt} that encodes risk-management heuristics. For instance, the LLM is explicitly instructed to prefer SELL or HOLD if the portfolio exposure exceeds 80\%, to avoid BUY recommendations if cash is depleted, and to flag high volatility coupled with negative sentiment as extreme risk.

    \section{Model Optimization and Training Setup}
    \textit{[Nota: Esta es su sección. Detallen el SFT Warm-up (1 época, 30 min) para enseñar el formato XML, la optimización Optuna TPE (3 trials, seed 42), y el entrenamiento GRPO (6 min). No olvidar explicar la Reward Function asimétrica $3\times3$ y cómo la aumentaron a -2.0 para el HOLD para romper el equilibrio de Nash.]}

    \section{Results and Validation}
    \textit{[Nota: Documenten aquí el contraste entre el Fast Mode y la Validación Completa (896 filas). Hablen del modelo ganador (exp\_20260506\_085649\_254046) y cómo se seleccionó manualmente por su Cumulative Return (90.67\%) a pesar de tener un MDD alto, contrastándolo con el modelo de mejor Fitness (1.974). Usen la tabla a continuación como base para sus datos.]}
    
    \begin{table}[H]
    \centering
    \begin{tabular}{@{}lcccc@{}}
    \toprule
    \textbf{Model ID} & \textbf{Fitness} & \textbf{Cumulative Return} & \textbf{Sharpe Ratio} & \textbf{MDD} \\ \midrule
    exp\_085649 (Selected) & 0.1069 & \textbf{90.67\%} & 2.9117 & 48.05\% \\
    exp\_111504 (Best Fit) & \textbf{1.9740} & 19.34\% & 1.9740 & \textbf{17.48\%} \\
    exp\_004508 (Best SR)  & -0.0759 & 89.76\% & \textbf{3.0408} & 51.17\% \\ \bottomrule
    \end{tabular}
    \caption{Full validation results comparing the selected model against other Optuna trials.}
    \end{table}

    \section{Conclusion and Future Work}
    In this Working Note, we presented a comprehensive RAG pipeline for financial decision-making in the CLEF 2026 FinMMEval Lab. By developing the \textit{finvert} module for precise sentiment extraction and a novel \textit{PromptBuilder} equipped with a tripartite temporal memory, we successfully bridged the gap between raw quantitative market data and qualitative LLM reasoning. 

    Our dual-strategy prompt orchestration demonstrated that providing an LLM with a structured, self-reflective memory of its past actions significantly enhances its ability to manage risk and make context-aware trading decisions. 

    Future work will focus on \textit{[agregar sus puntos sobre mejorar el split temporal (Walk-forward validation), balancear el dataset entre equity/crypto, y automatizar la selección por fitness en lugar de selección manual por CR.]}.

    \section*{Acknowledgments}
    We extend our gratitude to the Artificial Intelligence Club (GAIA) at ESCOM-IPN for their continuous support and computational resources. 

    \printbibliography

\end{document}